\documentclass[12pt]{article}

\usepackage{amsfonts}
\usepackage{amsmath}
\usepackage{amssymb}
\usepackage{amsthm}
\usepackage{color}
\usepackage[bottom]{footmisc}
\usepackage[a4paper, left=10mm, right=10mm, top=10mm, bottom=10mm]{geometry}
\usepackage{graphicx}
\usepackage{kotex}
\usepackage{listings}
\usepackage{setspace}
\usepackage{enumitem}
\usepackage{proof}
\usepackage{xparse}
\usepackage[utf8]{inputenc}
\usepackage{amssymb}
\usepackage{lmodern}
\usepackage{mathtools, nccmath}
\usepackage{mathtools}
\usepackage{indentfirst}
\usepackage{nopageno}
\allowdisplaybreaks
\sloppy

\newtheorem{theorem}{Theorem}
\newtheorem{fact}[theorem]{Fact}
\theoremstyle{definition}
\newtheorem{definition}[theorem]{Definition}

\DeclarePairedDelimiterX\set[1]\lbrace\rbrace{\def\mid{\;\delimsize\vert\;}#1}

\newenvironment{context}[1][]
{ \noindent \textbf{{#1}.} }
{ \hfill $ \dashv $ }

\newcommand{\powerset}{ \ensuremath{\mathcal{P}} }

\title{LA{(1)}LR{(0)}}
\author{임기정}

\begin{document}

\thispagestyle{empty}

\begin{definition}[$\mathrm{LRA}(G)$]
Let \(G \triangleq ( N , T , P , S )\) be a context-free grammar{.}
\begin{enumerate} \setlength{\itemsep}{0.2pt} \setlength{\parskip}{0.2pt}
\item[$\mathrm{(i)}$] $G^\prime := ( N^\prime , T^\prime , P^\prime , S^\prime ) \triangleq ( N \uplus \set*{ S^\prime } , T \uplus \set*{ \$ } , P \cup \set*{ [ S^\prime \to S \$ ] } , S^\prime )$.
\item[$\mathrm{(ii)}$] $V \triangleq N \uplus T$. $V^\prime \triangleq N^\prime \uplus T^\prime$.
\item[$\mathrm{(iii)}$] $\beta \xRightarrow[\mathbf{rm}]{}_{P^\prime} \gamma \xLeftrightarrow[]{\mathrm{def}} ( \beta , \gamma ) \in \set*{ ( \alpha A z , \alpha \omega z ) \mid [ A \to \omega ] \in P^\prime , \alpha \in V^{\prime \ast}, z \in T^{\prime \ast} }$.
\item[$\mathrm{(iv)}$] $I_{\mathrm{LR(0)}} \triangleq \set*{ [ A \to \alpha \cdot \beta ] \mid [ A \to \alpha \beta ] \in P^\prime }$.
\item[$\mathrm{(v)}$] $\mathrm{Cl} ( q ) =_\mu q \cup \set*{ [ A \to \varepsilon \cdot \omega ] \mid [ A \to \omega ] \in P^\prime , [ B \to \beta \cdot A \gamma ] \in \mathrm{Cl} ( q ) }$ for $ q \in \powerset ( I_{\mathrm{LR(0)}} ) $.
\item[$\mathrm{(vi)}$] $q_0 := \mathrm{Cl} ( \set*{ [ S^\prime \to \varepsilon \cdot S \$ ] } )$.
\item[$\mathrm{(vii)}$] $\mathrm{GOTO} ( q , X ) := \mathrm{Cl} ( \set*{ [ A \to \alpha X \cdot \beta ] \mid [ A \to \alpha \cdot X \beta ] \in q } )$ for $q \in \powerset ( I_{\mathrm{LR(0)}} )$ and $X \in V^\prime$.
\item[$\mathrm{(viii)}$] $Q \triangleq \mathrm{PT} \setminus \set*{ \emptyset }$ where $\mathrm{PT} =_\mu \set*{ q_0 } \cup \set*{ \mathrm{GOTO} ( q , X ) \mid q \in \mathrm{PT} , X \in V^\prime }$.
\item[$\mathrm{(ix)}$] $\varepsilon : p \to q \xLeftrightarrow[]{\mathrm{def}} p \in Q \land p = q$. $X \alpha : p \to q \xLeftrightarrow[]{\mathrm{def}} p \in Q \land \alpha : \mathrm{GOTO} ( p , X ) \to q$.
\item[$\mathrm{(x)}$] $\mathbf{Config} \triangleq \set*{ ( \alpha : p \to q , z ) \mid \alpha \in V^{\prime \ast}, p \in Q , q \in Q , z \in T^{\prime \ast} , \alpha : p \to q } $.
\item[$\mathrm{(xi)}$] $\delta ( q , X ) := \mathrm{GOTO} ( q , X )$ for $q \in Q$ and $X \in V^\prime$.
\item[$\mathrm{(xii)}$] $\mathbf{reduce} ( q , t ) := \set*{ [ A \to \omega ] \mid [ A \to \omega \cdot \varepsilon ] \in q }$ for $q \in Q$ and $t \in T^\prime$.
\item[$\mathrm{(xiii)}$] Let $\vdash$ be a binary relation on the set $\mathbf{Config}$ with two introduction rules:
\begin{center}
\( \infer[\mathrm{Shift}]{ ( \alpha : q \to q^\prime , t z ) \vdash ( \alpha t : q \to q^{\prime \prime} , z ) }{ q^{\prime \prime} = \delta ( q^\prime , t ) } \) 

\hfil

\( \infer[\mathrm{Reduce} ( A \to \omega )]{ ( \alpha \omega : q \to q^\prime , t z ) \vdash ( \alpha A : q \to q^{\prime \prime} , t z ) }{ \alpha : q \to p & \omega : p \to q^\prime & [ A \to \omega ] \in \mathbf{reduce} ( q^\prime , t ) & q^{\prime\prime} = \delta (p , A) } \)
\end{center}
\item[$\mathrm{(xiv)}$] $q_f := \delta ( \delta ( q_0 , S ) , \$ )$.
\item[$\mathrm{(xv)}$] The language accepted by $\mathrm{LRA}(G)$ is defined as follows: \[L ( \mathrm{LRA} ( G ) ) := \set*{ z \in { T }^\ast \mid ( \varepsilon : q_0 \to q_0 , z \$ ) \vdash^\ast ( S \$ : q_0 \to q_f , \varepsilon ) } .\]
\end{enumerate}
\end{definition}

\begin{fact}
It is known that $L ( G ) = L ( \mathrm{LRA} ( G ) )$, where \(L(G) \triangleq \set*{ z \in { T }^\ast \mid S \Rightarrow_{P}^\ast z }\).
\end{fact}

\begin{proof}
For a configuration \(c=(\alpha:p\to q,u)\in\mathbf{Config}\) write
\(\mathrm{Y}(c):=\alpha u\in V'^*\) for its \emph{yield}.

\paragraph{Viable-prefix lemma.}
For \(\eta\in V'^*\):
\begin{itemize}
\item[$(\mathrm V1)$] \((\exists q\in Q)\ \eta:q_0\to q\) iff \(\eta\) is a viable
prefix of \(G'\);
\item[$(\mathrm V2)$] if \(\eta=\eta'\omega\) with \(\eta':q_0\to p\),
\(\omega:p\to q\), and this occurrence of \(\omega\) is the handle of
\([A\to\omega]\in P'\) (i.e.
\(S'\xRightarrow[\mathbf{rm}]{}_{P'}^*\eta' A y\xRightarrow[\mathbf{rm}]{}_{P'}\eta'\omega y\),
\(y\in T'^*\)), then \([A\to\omega\cdot\varepsilon]\in q\), so
\([A\to\omega]\in\mathbf{reduce}(q,t)\) for the first symbol \(t\) of \(y\).
\end{itemize}
Both are the standard correctness of the characteristic LR(0) automaton and
follow by induction on the clauses \(\mathrm{Cl}\), \(\mathrm{GOTO}\) of
Definition 1(v),(vii) and the path relation (ix).

\paragraph{Step invariant.}
For \(c\vdash c'\),
\begin{equation}
\mathrm{Y}(c')=\mathrm{Y}(c)\ \ (\text{Shift}),\qquad
\mathrm{Y}(c')\xRightarrow[\mathbf{rm}]{}_{P'}\mathrm{Y}(c)\ \ (\text{Reduce}).
\tag{$\ast$}
\end{equation}
Indeed a Shift \((\alpha:q\to q',tz)\vdash(\alpha t:q\to q'',z)\) has
\(\mathrm{Y}=\alpha tz\) on both sides, and a Reduce
\((\alpha\omega:q\to q',tz)\vdash(\alpha A:q\to q'',tz)\), \([A\to\omega]\in P'\),
has \(\mathrm{Y}(c)=\alpha\omega tz\), \(\mathrm{Y}(c')=\alpha A tz\) with
\(\alpha\in V'^*\), \(tz\in T'^*\), so
\(\alpha A tz\xRightarrow[\mathbf{rm}]{}_{P'}\alpha\omega tz\) by Definition 1(iii).
Iterating \((\ast)\) along \(c\vdash^* c'\),
\begin{equation}
c\vdash^* c'\ \Longrightarrow\ \mathrm{Y}(c')\xRightarrow[\mathbf{rm}]{}_{P'}^*\mathrm{Y}(c).
\tag{$\ast\ast$}
\end{equation}

\paragraph{Soundness, \(L(\mathrm{LRA}(G))\subseteq L(G)\).}
Let \(z\in L(\mathrm{LRA}(G))\), so
\((\varepsilon:q_0\to q_0,z\$)\vdash^*(S\$:q_0\to q_f,\varepsilon)\). The endpoints
have yields \(z\$\) and \(S\$\), so \((\ast\ast)\) gives
\[
S\$\ \xRightarrow[\mathbf{rm}]{}_{P'}^*\ z\$.
\]
Since \(S'\) occurs on no right-hand side of \(P'\), \([S'\to S\$]\) is applied at
no step of a derivation from \(S\$\); thus every step uses a production of \(P\).
As \(\$\in T'\setminus T\) is never rewritten and never matched by a production of
\(P\), it rides unchanged at the right end, and deleting it gives
\(S\Rightarrow_P^* z\), i.e. \(z\in L(G)\).

\paragraph{Completeness, \(L(G)\subseteq L(\mathrm{LRA}(G))\).}
Let \(z\in L(G)\) and fix a rightmost derivation
\[
S\$=\gamma_0\xRightarrow[\mathbf{rm}]{}_{P'}\gamma_1\xRightarrow[\mathbf{rm}]{}_{P'}\cdots\xRightarrow[\mathbf{rm}]{}_{P'}\gamma_N=z\$,
\]
which exists because \(S\Rightarrow_P^* z\). Write the \(k\)-th step as
\[
\gamma_k=\alpha_k A_k y_k\ \xRightarrow[\mathbf{rm}]{}_{P'}\ \alpha_k\omega_k y_k=\gamma_{k+1},
\qquad [A_k\to\omega_k]\in P',\ \alpha_k\in V'^*,\ y_k\in T'^*,
\]
where \(y_k\) ends in \(\$\), hence is nonempty with first symbol \(t_k\). The
suffix \(\omega_k\) of the viable prefix \(\alpha_k\omega_k\) is the handle of
\([A_k\to\omega_k]\), so by \((\mathrm V2)\) there are \(p_k,q_k\in Q\) with
\[
\alpha_k:q_0\to p_k,\quad \omega_k:p_k\to q_k,\quad
[A_k\to\omega_k\cdot\varepsilon]\in q_k,\quad
[A_k\to\omega_k]\in\mathbf{reduce}(q_k,t_k).
\]
Read the derivation from \(\gamma_N=z\$\) down to \(\gamma_0=S\$\): starting at
\((\varepsilon:q_0\to q_0,z\$)\), for \(k=N-1,\dots,0\) Shift the terminals lying
between the current stack and the handle \(\omega_k\) (each legal by the Shift
rule, as \(\delta\) is defined along the viable prefix \(\alpha_k\omega_k\)) until
the configuration is \((\alpha_k\omega_k:q_0\to q_k,y_k)\), then apply
\(\mathrm{Reduce}(A_k\to\omega_k)\), obtaining by the line above and
Definition 1(xiii)
\[
(\alpha_k\omega_k:q_0\to q_k,\,y_k)\ \vdash\ (\alpha_k A_k:q_0\to\delta(p_k,A_k),\,y_k),
\]
the configuration associated with \(\gamma_k\). Each iteration inverts one
\(\xRightarrow[\mathbf{rm}]{}_{P'}\) step; after all \(N\) of them and a final Shift of the
endmarker \(\$\) the input is exhausted and the stack carries \(S\$\):
\[
(\varepsilon:q_0\to q_0,z\$)\ \vdash^*\ (S\$:q_0\to q_f,\varepsilon),
\qquad q_f=\delta(\delta(q_0,S),\$),
\]
so \(z\in L(\mathrm{LRA}(G))\). With soundness, \(L(G)=L(\mathrm{LRA}(G))\).
\end{proof}


\begin{theorem}
Define $\mathbf{LA} : \set*{ ( q , [ A \to \omega ] ) \mid q \in Q , [ A \to \omega \cdot \varepsilon ] \in q } \to \powerset ( T^\prime )$ by
\begin{align}
\mathbf{LA} ( q , [ A \to \omega ] ) := \set*{ t \in T^\prime \mid S^\prime \xRightarrow[\mathbf{rm}]{}_{P^\prime}^\ast \alpha A t z , \alpha \omega : q_0 \to q , \alpha \in V^{\prime \ast}, z \in T^{\prime \ast} } .
\end{align}
Then, overriding $\mathbf{reduce} : Q \times T^\prime \to \powerset ( P^\prime )$ of the $\mathrm{LR(0)}$ parser $\mathrm{LRA} ( G )$ with
\begin{align*}
( q , t ) \mapsto \set*{ [ A \to \omega ] \mid [ A \to \omega \cdot \varepsilon ] \in q , t \in \mathbf{LA} ( q , [ A \to \omega ] ) }
\end{align*}
yields an $\mathrm{LALR(1)}$ parser if there are no conflicts, which accepts the same language.
\end{theorem}

\begin{proof}
Let \(\vdash_{\mathbf{LA}}\) be \(\vdash\) with \(\mathbf{reduce}\) replaced by
\[
\mathbf{reduce}_{\mathbf{LA}}(q,t):=\set*{[A\to\omega]\mid [A\to\omega\cdot\varepsilon]\in q,\ t\in\mathbf{LA}(q,[A\to\omega])},
\]
and let \(L(\vdash_{\mathbf{LA}})\) be the language it accepts.

\(\boldsymbol{(\subseteq)}\) Since
\(\mathbf{reduce}_{\mathbf{LA}}(q,t)\subseteq\mathbf{reduce}(q,t)\) for every
\((q,t)\), each Reduce instance of \(\vdash_{\mathbf{LA}}\) is one of \(\vdash\),
and the Shift rule is shared; hence \(\vdash_{\mathbf{LA}}\subseteq\vdash\),
\(\vdash_{\mathbf{LA}}^*\subseteq\vdash^*\), and
\[
L(\vdash_{\mathbf{LA}})\subseteq L(\mathrm{LRA}(G))\overset{\text{Fact 2}}{=}L(G).
\]

\(\boldsymbol{(\supseteq)}\) Let \(z\in L(G)\). The completeness construction of
Fact 2 yields a computation
\[
(\varepsilon:q_0\to q_0,z\$)\ \vdash^*\ (S\$:q_0\to q_f,\varepsilon)
\]
each of whose Reduce steps inverts a rightmost step
\(\alpha A tz'\xRightarrow[\mathbf{rm}]{}_{P'}\alpha\omega tz'\) and is taken at a
configuration \((\alpha\omega:q_0\to q,tz')\) with
\([A\to\omega\cdot\varepsilon]\in q\) and \(\alpha\omega:q_0\to q\). As
\(S'\xRightarrow[\mathbf{rm}]{}_{P'}S\$\xRightarrow[\mathbf{rm}]{}_{P'}^*\alpha A tz'\), the
witnesses \(\alpha,z'\) and the path \(\alpha\omega:q_0\to q\) meet the defining
condition (1) of \(\mathbf{LA}\), so
\[
t\in\mathbf{LA}(q,[A\to\omega]),\qquad\text{i.e.}\qquad
[A\to\omega]\in\mathbf{reduce}_{\mathbf{LA}}(q,t).
\]
Thus every step is also a \(\vdash_{\mathbf{LA}}\)-step and the same computation
accepts \(z\); hence \(L(G)\subseteq L(\vdash_{\mathbf{LA}})\). Combining,
\(L(\vdash_{\mathbf{LA}})=L(G)\).

For the table, each pair \((q,t)\) carries the shift to \(\delta(q,t)\) when
\(\delta(q,t)\ne\emptyset\) and the reduces \([A\to\omega]\) with
\([A\to\omega\cdot\varepsilon]\in q\) and \(t\in\mathbf{LA}(q,[A\to\omega])\)~---
exactly the LR(0) cores annotated with the LALR(1) lookahead sets \(\mathbf{LA}\).
If no \((q,t)\) enables two distinct actions~--- no shift/reduce
(\(\delta(q,t)\ne\emptyset\) together with an enabled reduce) and no
reduce/reduce (two distinct enabled reduces)~--- then \(\vdash_{\mathbf{LA}}\) is
deterministic with one symbol of lookahead, i.e. an \(\mathrm{LALR(1)}\) parser.
\end{proof}


\begin{theorem}
Letting $ R \mathbin{!} x \triangleq \set*{ y \mid ( x , y ) \in R } $, define relations $\mathbf{Read}$ and $\mathbf{Follow}$ from the set \[ \set*{ ( p , A ) \mid p \in Q , A \in N^\prime , \delta ( p , A ) \ne \emptyset } \] to the set $T^\prime$ inductively as follows:
\[
\begin{matrix}
\infer[\mathrm{DR}]{ t \in \mathbf{Read} \mathbin{!} ( p , A ) }{ \delta ( \delta ( p , A ) , t ) \ne \emptyset } & & \infer[\mathrm{reads}]{ \mathbf{Read} \mathbin{!} ( r , C ) \subseteq \mathbf{Read} \mathbin{!} ( p , A ) }{ \delta ( p , A ) = r & C \Rightarrow_{P^\prime}^\ast \varepsilon } \\
& & \\
\infer[\mathrm{Read}]{ t \in \mathbf{Follow} \mathbin{!} ( p , A ) }{ t \in \mathbf{Read} \mathbin{!} ( p , A ) } & & \infer[\mathrm{includes}]{ \mathbf{Follow} \mathbin{!} ( p^\prime , B ) \subseteq \mathbf{Follow} \mathbin{!} ( p , A ) }{ [ B \to \beta \cdot A \gamma ] \in p & \beta : p^\prime \to p & \gamma \Rightarrow_{P^\prime}^\ast \varepsilon }
\end{matrix}
\]
Then, whenever $q \in Q$ and $[ A \to \omega \cdot \varepsilon ] \in q$, we can compute $\mathbf{LA} (q , [ A \to \omega ])$ by
\begin{align}
\mathbf{LA} ( q , [ A \to \omega ] ) = \set*{ t \in T^\prime \mid p \in Q , \omega : p \to q , \delta ( p , A ) \ne \emptyset , t \in \mathbf{Follow} \mathbin{!} ( p , A ) } .
\end{align}
\end{theorem}

\begin{proof}
Let
\[
D=\set*{(p,A)\mid p\in Q,\ A\in N',\ \delta(p,A)\ne\emptyset}.
\]
For the proof, define two semantic sets on \(D\).  First, let
\[
\mathsf{Read}(p,A)=
\set*{t\in T' \mid
\delta(p,A)=r,\ \gamma\Rightarrow_{P'}^*\varepsilon,\
\gamma t:r\to s\text{ for some }\gamma\in N'^*,\ s\in Q } .
\]
Thus \(t\in\mathsf{Read}(p,A)\) exactly when, after the transition on \(A\), one
can reach a terminal transition on \(t\) through zero or more nullable
nonterminal transitions.  Second, let
\[
\mathsf{Follow}(p,A)=
\set*{t\in T' \mid
S'\xRightarrow[\mathbf{rm}]{}_{P'}^*\alpha A tz,\
\alpha:q_0\to p,\ \alpha\in V'^*,\ z\in T'^* } .
\]
This is the semantic follow set of the nonterminal transition \((p,A)\).

\paragraph{Computation of \(\mathsf{Read}\).}
We prove \(\mathbf{Read}\mathbin{!}(p,A)=\mathsf{Read}(p,A)\) on \(D\) by showing
both are the least fixed point of \(\{\mathrm{DR},\mathrm{reads}\}\). For
\(t\in\mathsf{Read}(p,A)\) fix a witness \(\gamma\in N'^*\) with
\(\gamma\Rightarrow_{P'}^*\varepsilon\) and \(\gamma t:\delta(p,A)\to s\), and
induct on \(|\gamma|\):
\[
\begin{aligned}
|\gamma|=0:\quad& t:\delta(p,A)\to s
\ \overset{\text{1(ix)}}{\Longleftrightarrow}\ s=\delta(\delta(p,A),t)\ne\emptyset
\ \Longleftrightarrow\ t\in\mathrm{DR}(p,A);\\
\gamma=C\gamma'\ (C\in N'):\quad& C\Rightarrow_{P'}^*\varepsilon,\
\gamma'\Rightarrow_{P'}^*\varepsilon,\
\gamma' t:\delta(\delta(p,A),C)\to s,
\end{aligned}
\]
so in the second case, with \(r=\delta(p,A)\), we have \((r,C)\in D\),
\((p,A)\leadsto_{\mathrm{reads}}(r,C)\), and \(t\in\mathsf{Read}(r,C)\) via the
shorter witness \(\gamma'\); the induction hypothesis gives
\(t\in\mathbf{Read}\mathbin{!}(r,C)\), whence \(\mathrm{reads}\) yields
\(t\in\mathbf{Read}\mathbin{!}(p,A)\). The reverse inclusion reads the same
decomposition forwards (\(\mathrm{DR}\) realizes \(\gamma=\varepsilon\); each
\(\mathrm{reads}\) edge prepends one nullable \(C\in N'\)). As \(D,T'\) are finite
and both clauses are monotone, the two least fixed points coincide:
\[
\mathbf{Read}\mathbin{!}(p,A)=\mathsf{Read}(p,A)\qquad((p,A)\in D),
\]
with \(\mathbf{Read}\mathbin{!}(r,C)=\emptyset\) when \((r,C)\notin D\), making
the \(\mathrm{reads}\) clause vacuous there.

\paragraph{Computation of \(\mathsf{Follow}\).}
We show \(\mathsf{Follow}\) is the least solution of the clauses \(\mathrm{Read}\)
and \(\mathrm{includes}\), whence \(\mathbf{Follow}\mathbin{!}=\mathsf{Follow}\) on
\(D\).

\emph{Soundness of the clauses for \(\mathsf{Follow}\).}
\begin{itemize}
\item[$(\mathrm{Read})$] \(\mathsf{Read}(p,A)\subseteq\mathsf{Follow}(p,A)\). Take
\(t\in\mathsf{Read}(p,A)\) with witness \(\gamma\in N'^*\),
\(\gamma\Rightarrow_{P'}^*\varepsilon\), \(\gamma t:\delta(p,A)\to s\). Choose
\(\alpha:q_0\to p\) (possible since \(p\in Q\) is accessible); then
\(\alpha A\gamma t:q_0\to s\) is a viable prefix, so some rightmost derivation
reaches \(\alpha A\gamma t z\) (\(z\in T'^*\)). As \(\gamma\) stands immediately
left of the terminals \(tz\), it is rightmost-erasable, giving
\[
S'\xRightarrow[\mathbf{rm}]{}_{P'}^*\alpha A\gamma t z\xRightarrow[\mathbf{rm}]{}_{P'}^*\alpha A t z,
\qquad\alpha:q_0\to p,
\]
i.e. \(t\in\mathsf{Follow}(p,A)\).
\item[$(\mathrm{includes})$] If \([B\to\beta\cdot A\gamma]\in p\),
\(\beta:p'\to p\), \(\gamma\Rightarrow_{P'}^*\varepsilon\), then
\(\mathsf{Follow}(p',B)\subseteq\mathsf{Follow}(p,A)\). For
\(t\in\mathsf{Follow}(p',B)\), say \(S'\xRightarrow[\mathbf{rm}]{}_{P'}^*\alpha' B t z\)
with \(\alpha':q_0\to p'\), expand \(B\) and erase \(\gamma\) rightmost:
\[
\alpha' B t z\ \xRightarrow[\mathbf{rm}]{}_{P'}\ \alpha'\beta A\gamma t z\ \xRightarrow[\mathbf{rm}]{}_{P'}^*\ \alpha'\beta A t z,
\qquad \alpha'\beta:q_0\to p,
\]
where \(\alpha'\beta:q_0\to p\) follows from \(\alpha':q_0\to p'\) and
\(\beta:p'\to p\). With \(\alpha:=\alpha'\beta\) this gives
\(t\in\mathsf{Follow}(p,A)\).
\end{itemize}

\emph{Minimality.} Let \(t\in\mathsf{Follow}(p,A)\), witnessed by
\(S'\xRightarrow[\mathbf{rm}]{}_{P'}^*\alpha A t z\), \(\alpha:q_0\to p\). Let
\([B\to\beta A\gamma]\) be the production creating this occurrence of \(A\), so the
derivation factors through
\[
S'\xRightarrow[\mathbf{rm}]{}_{P'}^*\alpha' B t' z'\ \xRightarrow[\mathbf{rm}]{}_{P'}\ \alpha'\beta A\gamma t' z'\ \xRightarrow[\mathbf{rm}]{}_{P'}^*\ \alpha A t z,
\quad \alpha=\alpha'\beta,\ [B\to\beta\cdot A\gamma]\in p,\ \beta:p'\to p .
\]
Splitting on the residual \(\gamma\):
\[
\begin{cases}
\gamma=\gamma_0\,t\cdots,\ \gamma_0\Rightarrow_{P'}^*\varepsilon
& \Rightarrow\ \gamma_0 t:\delta(p,A)\to(\cdot),\ \text{so } t\in\mathsf{Read}(p,A)\ \ (\mathrm{Read}),\\[2pt]
\gamma\Rightarrow_{P'}^*\varepsilon
& \Rightarrow\ t\in\mathsf{Follow}(p',B)\ \text{via } (p,A)\leadsto_{\mathrm{includes}}(p',B).
\end{cases}
\]
The second case strictly shortens the witnessing rightmost derivation, so
iterating it terminates in the first case after finitely many steps; thus every
\(t\in\mathsf{Follow}(p,A)\) is generated by the clauses. Hence \(\mathsf{Follow}\)
is their least solution and
\[
\mathbf{Follow}\mathbin{!}(p,A)=\mathsf{Follow}(p,A)\qquad((p,A)\in D).
\]

It remains to translate this semantic follow set into the lookahead set of a
completed item.  Let \(q\in Q\) and \([A\to\omega\cdot\varepsilon]\in q\).  Then
\[
\begin{aligned}
&t\in\mathbf{LA}(q,[A\to\omega]) \\
&\quad\Longleftrightarrow
\exists\alpha,z\;.
S'\xRightarrow[\mathbf{rm}]{}_{P'}^*\alpha A tz
\land \alpha\omega:q_0\to q \\
&\quad\Longleftrightarrow
\exists p\in Q,
\omega:p\to q \land \delta(p,A)\ne\emptyset
\land t\in\mathsf{Follow}(p,A) \\
&\quad\Longleftrightarrow
\exists p\in Q,
\omega:p\to q \land \delta(p,A)\ne\emptyset
\land t\in\mathbf{Follow}\mathbin{!}(p,A).
\end{aligned}
\]
The middle equivalence just cuts the path \(\alpha\omega:q_0\to q\) at the point
\(p\) reached after \(\alpha\); conversely, concatenating \(\alpha:q_0\to p\) and
\(\omega:p\to q\) recovers \(\alpha\omega:q_0\to q\).  The side condition
\(\delta(p,A)\ne\emptyset\) on the second line is automatic when \(A\ne S'\):
from \(\omega:p\to q\) and \([A\to\omega\cdot\varepsilon]\in q\) one traces the dot
back along \(\omega\) to obtain \([A\to\varepsilon\cdot\omega]\in p\), and this
closure item is licensed by some \([B\to\beta\cdot A\delta]\in p\), whence
\(\delta(p,A)\ne\emptyset\).  For \(A=S'\) both sides are empty: no item has a dot
before \(S'\), so \(\delta(p,S')=\emptyset\) for every \(p\), and no rightmost
derivation has the form \(S'\xRightarrow[\mathbf{rm}]{}_{P'}^*\alpha S' tz\) since
\(S'\) occurs on no right-hand side; this matches the parser accepting by reaching
\((S\$:q_0\to q_f,\varepsilon)\) rather than by reducing \([S'\to S\$]\).  This is
precisely the formula claimed in the theorem.
\end{proof}


\appendix
\section*{Appendix A. Digraph algorithm and implementation specification}

This appendix gives an executable specification for computing the sets used in
Theorem 4.  The construction is the usual digraph view of LALR(1) lookahead
computation: direct seed sets are placed on nodes, dependency edges describe set
inclusions, and the least fixed point is obtained by collapsing strongly
connected components and propagating set unions along the resulting DAG
\cite{DeRemerPennello1982,Tarjan1972,Gallier2008}.

\subsection*{A.1. General digraph fixed-point problem}
Let \(X\) be a finite set of nodes, let \(E\subseteq X\times X\) be a directed
edge relation, and let \(\mathrm{Seed}:X\to\powerset(T')\).  We use the edge
orientation
\[
  x\leadsto y \quad\text{to mean}\quad F(y)\subseteq F(x).
\]
Thus the desired solution is the least function \(F:X\to\powerset(T')\)
satisfying
\begin{align}
F(x)=\mathrm{Seed}(x)\cup\bigcup_{x\leadsto y}F(y). \tag{DG}
\end{align}

\paragraph{Algorithm \(\mathrm{DIGRAPH}(X,E,\mathrm{Seed})\).}
\begin{enumerate}[leftmargin=7mm]
\item Compute the strongly connected components of \((X,E)\), for example by
Tarjan's depth-first-search algorithm.
\item Build the condensation DAG \(\mathcal C\).  Its nodes are SCCs, and
\(C\leadsto D\) is an edge when some \(x\in C\) and \(y\in D\) satisfy
\(x\leadsto y\) and \(C\ne D\).
\item Initialize
\[
  \mathrm{Val}(C)\leftarrow\bigcup_{x\in C}\mathrm{Seed}(x)
\]
for every SCC \(C\).
\item Process the SCCs in reverse topological order with respect to
\(\leadsto\), so that every successor \(D\) of \(C\) has already been processed,
and put
\[
  \mathrm{Val}(C) \leftarrow \mathrm{Val}(C)\cup
  \bigcup_{C\leadsto D}\mathrm{Val}(D).
\]
\item Return \(F(x):=\mathrm{Val}(C_x)\), where \(C_x\) is the SCC containing
\(x\).
\end{enumerate}

\paragraph{Correctness condition.}
For every SCC \(C\), after step 4,
\[
  \mathrm{Val}(C)=\bigcup\{\mathrm{Seed}(y)\mid
  \text{some }x\in C\text{ reaches }y\text{ in }(X,E)\}.
\]
Therefore the returned \(F\) is exactly the least solution of \((\mathrm{DG})\).
This is the only property of the graph algorithm needed by Theorem 4.

\paragraph{Complexity.}
Let \(b=\lceil |T'|/w\rceil\), where \(w\) is the machine word size used for
bitsets.  SCC construction is \(O(|X|+|E|)\), and the set propagation is
\(O((|X|+|E|)b)\).  A simple work-list implementation that repeatedly applies
\(F(x):=F(x)\cup F(y)\) for \(x\leadsto y\) is also correct, but the SCC version
avoids repeated propagation inside cycles.

\subsection*{A.2. Instance for \(\mathbf{Read}\)}
Let
\[
D=\set*{(p,A)\mid p\in Q,\ A\in N',\ \delta(p,A)\ne\emptyset}.
\]
First compute nullable nonterminals by the least fixed point
\[
  \mathrm{Null}(A) \Longleftrightarrow
  \exists [A\to X_1\cdots X_k]\in P'\text{ such that every }X_i
  \text{ is nullable},
\]
where terminals are never nullable and \(k=0\) is allowed.  Extend this to
strings by \(\mathrm{Null}(X_1\cdots X_k)\) iff every \(X_i\) is nullable.

For each node \(x=(p,A)\in D\), define the direct-read seed
\[
  \mathrm{DR}(p,A)=\set*{t\in T'\mid \delta(\delta(p,A),t)\ne\emptyset}.
\]
Define the reads dependency graph \(E_R\subseteq D\times D\) by
\[
  (p,A)\leadsto_R(r,C)
  \quad\Longleftrightarrow\quad
  \delta(p,A)=r \land C\in N' \land
  \mathrm{Null}(C) \land
  (r,C)\in D.
\]
The last condition merely makes explicit the domain of Theorem 4.  If
\((r,C)\notin D\), then \(\mathbf{Read}\mathbin{!}(r,C)=\emptyset\), so no edge is
needed.

Then compute
\[
  \mathbf{Read}\mathbin{!}(p,A)
  :=\mathrm{DIGRAPH}(D,E_R,\mathrm{DR})(p,A).
\]
This is the least relation satisfying the \(\mathrm{DR}\) and
\(\mathrm{reads}\) rules of Theorem 4.

\subsection*{A.3. Instance for \(\mathbf{Follow}\)}
The seed for follow propagation is the already computed read set:
\[
  \mathrm{Seed}_F(p,A):=\mathbf{Read}\mathbin{!}(p,A).
\]
Define the includes dependency graph \(E_I\subseteq D\times D\) by
\[
(p,A)\leadsto_I(p',B)
\quad\Longleftrightarrow\quad
[B\to\beta\cdot A\gamma]\in p \land
\beta:p'\to p \land
\mathrm{Null}(\gamma) \land
(p',B)\in D.
\]
The edge orientation again means that the right-hand node contributes to the
left-hand node:
\[
  \mathbf{Follow}\mathbin{!}(p',B)
  \subseteq
  \mathbf{Follow}\mathbin{!}(p,A).
\]
Then compute
\[
  \mathbf{Follow}\mathbin{!}(p,A)
  :=\mathrm{DIGRAPH}(D,E_I,\mathrm{Seed}_F)(p,A).
\]
This is the least relation satisfying the \(\mathrm{Read}\) and
\(\mathrm{includes}\) rules of Theorem 4.

\subsection*{A.4. Lookback relation and lookahead construction}
For every completed item \([A\to\omega\cdot\varepsilon]\in q\), define its
lookback set by
\[
  \mathrm{LB}(q,A\to\omega)
  :=\set*{(p,A)\in D\mid \omega:p\to q}.
\]
The lookahead set is then computed by
\[
  \mathbf{LA}(q,[A\to\omega])
  :=\bigcup_{(p,A)\in\mathrm{LB}(q,A\to\omega)}
  \mathbf{Follow}\mathbin{!}(p,A).
\]
This is the formula of Theorem 4 written as an implementation step.

\paragraph{Implementation note.}
The query \(\omega:p\to q\) can be implemented by following \(\delta\)-edges from
\(p\) along \(\omega\).  The query \(\beta:p'\to p\) in \(E_I\) can be implemented
by caching reverse-path queries keyed by \((\beta,p)\), or by storing predecessor
edges during construction of the LR(0) automaton.

\subsection*{A.5. Parser-table and conflict specification}
After \(\mathbf{LA}\) has been computed, construct actions as follows.
\begin{enumerate}[leftmargin=7mm]
\item If \(\delta(q,t)\ne\emptyset\) for \(t\in T'\), add the shift action
\(q\xrightarrow{t}\delta(q,t)\).
\item For each \([A\to\omega\cdot\varepsilon]\in q\) and each
\(t\in\mathbf{LA}(q,[A\to\omega])\), add the reduce action
\(A\to\omega\) under lookahead \(t\).
\item Report a shift/reduce conflict when both a shift action and a reduce
action are present for the same pair \((q,t)\).
\item Report a reduce/reduce conflict when two distinct reduce productions are
present for the same pair \((q,t)\).
\end{enumerate}
The accepting condition remains the one in Definition 1: the parser accepts when
it reaches \((S\$:q_0\to q_f,\varepsilon)\).  Equivalently, in a conventional
ACTION/GOTO table, \(q_f\) may be treated as the accepting state after the input
marker has been consumed.

\subsection*{A.6. Summary specification}
The complete computation required by Theorem 4 is:
\begin{enumerate}[leftmargin=7mm]
\item Build \(G'\), \(Q\), and \(\delta\) as in Definition 1.
\item Compute nullable nonterminals and nullable strings.
\item Build \(D\), \(\mathrm{DR}\), and \(E_R\); compute \(\mathbf{Read}\) by
\(\mathrm{DIGRAPH}(D,E_R,\mathrm{DR})\).
\item Build \(E_I\); compute \(\mathbf{Follow}\) by
\(\mathrm{DIGRAPH}(D,E_I,\mathbf{Read})\).
\item Build \(\mathrm{LB}\) and compute each \(\mathbf{LA}(q,[A\to\omega])\) by
unioning the appropriate \(\mathbf{Follow}\)-sets.
\item Construct the guarded reduce actions and check conflicts.
\end{enumerate}
The postcondition is that the computed \(\mathbf{Read}\), \(\mathbf{Follow}\),
and \(\mathbf{LA}\) are the least sets satisfying Theorem 4, and hence the
lookahead-restricted parser of Theorem 3 uses exactly the LALR(1) lookahead sets
determined by the LR(0) automaton.

\begin{thebibliography}{9}
\bibitem{DeRemerPennello1982}
F. DeRemer and T. J. Pennello.
\newblock Efficient computation of LALR(1) look-ahead sets.
\newblock \emph{ACM Transactions on Programming Languages and Systems},
4(4):615--649, 1982.

\bibitem{Tarjan1972}
R. E. Tarjan.
\newblock Depth-first search and linear graph algorithms.
\newblock \emph{SIAM Journal on Computing}, 1(2):146--160, 1972.

\bibitem{Gallier2008}
J. H. Gallier.
\newblock A survey of LR-parsing methods: The graph method for computing fixed
points and LALR(1) lookahead sets.
\newblock CIS 511 lecture notes, University of Pennsylvania, 2008.
\end{thebibliography}

\end{document}
